% G3-6 Tahapan eksperimen.
% Sumber: diolah penulis.
% Tidak memuat hasil pemilihan tahap mana pun; rincian setiap tahap ada di
% narasi Subbab Skenario Pengujian.
\begin{tikzpicture}[x=1mm,y=-1mm,font=\footnotesize,
    tahap/.style={draw,rectangle,text width=88mm,align=left,
      inner sep=3pt,minimum height=11mm},
    label/.style={draw,rectangle,fill=black!10,text width=26mm,
      align=center,inner sep=3pt,minimum height=11mm,font=\small},
    panah/.style={-{Stealth[length=2.2mm]}}]

  % Tinggi kotak 11mm dan jarak antarbaris 17mm, seragam untuk semua tahap.
  \node[label] (l1) at (14,0) {Tahap 1\\\textit{time frame}};
  \node[tahap] (t1) at (74,0)
    {9 pasangan \textit{dataset} dan horizon $\times$ 4 model.\\
     Hasil: \textit{dataset} dan horizon untuk tahap berikutnya.};

  \node[label] (l2) at (14,17) {Tahap 2\\metode optimasi};
  \node[tahap] (t2) at (74,17)
    {7 metode optimasi $\times$ 4 model.\\
     Hasil: metode optimasi untuk tahap berikutnya.};

  \node[label] (l3) at (14,34) {Tahap 3\\rasio \textit{split}};
  \node[tahap] (t3) at (74,34)
    {3 rasio \textit{split} $\times$ 4 model.\\
     Tanpa pemilihan; semua hasil dilaporkan.};

  \node[label] (l4) at (14,51) {Tahap 4\\teknik tambahan};
  \node[tahap] (t4) at (74,51)
    {4 teknik tambahan $\times$ 4 model.\\
     Tanpa pemilihan; semua hasil dilaporkan.};

  \foreach \a/\b in {1/2,2/3,3/4}{
    \draw[panah] (t\a.south) -- (t\b.north);
  }
\end{tikzpicture}
